\chapter{Ejercicios}
\label{cap:ejercicios}

Una guía se aprende haciendo. Estos ejercicios están ordenados de menos a más y tocan, cada
uno, una parte distinta del proyecto. Para todos:

\begin{itemize}
  \item Crea una rama (\cod{git switch -c practica/ejercicio-N}) y trabaja en ella; los
  ejercicios son para tirar.
  \item Usa siempre un \cod{HOME} temporal y \cod{MUSICPLAYER_SILENT=1}
  (capítulo~\ref{cap:depurar}).
  \item Si te atascas, busca en el proyecto algo que ya haga algo parecido: casi todo tiene un
  vecino que puedes imitar. Al final del capítulo hay pistas para cada uno, y después soluciones
  resumidas.
\end{itemize}

\section{Nivel 1: solo QML}

\begin{ejercicio}[1.1 Tiempo que falta]
En la barra de reproducción (\archivo{PlayerBar.qml}) muestra el tiempo que \emph{falta}
(\cod{-1:23}) en vez de la duración total. Haz la prueba con una captura.
\end{ejercicio}

\begin{ejercicio}[1.2 Hacerlo una opción]
Convierte lo anterior en una opción de Opciones («Tiempo restante», sí o no). Sigue la receta
de la capítulo~\ref{cap:recetas}. Comprueba que \cod{docs.py check} te exige documentarla en la
ficha del área.
\end{ejercicio}

\begin{ejercicio}[1.3 Un color nuevo]
Añade a las skins un color \cod{tiempo} para ese texto. Pon un valor distinto en la skin Neón.
Haz capturas con dos skins.
\end{ejercicio}

\section{Nivel 2: Rust sin puente}

\begin{ejercicio}[2.1 Una sonda]
Añade la sonda \cod{--probe artist NOMBRE}, que busque y muestre solo el primer artista
(título y \cod{browseId}). Úsala sin abrir ninguna ventana.
\end{ejercicio}

\begin{ejercicio}[2.2 Una prueba]
Escribe una prueba para \cod{seconds("3:27")} en \archivo{ytm.rs}: debe dar 207. Añade
también \cod{"1:02:03"} (3\,723) y la cadena vacía (0). Ejecuta
\cod{cargo test --release}.
\end{ejercicio}

\begin{ejercicio}[2.3 Un parseador]
En \archivo{lyrics.rs}, la función \cod{parse} lee el formato LRC. Escribe una prueba para una
línea con \emph{dos} marcas de tiempo (\cod{[00:10.00][01:30.00]Estribillo}): debe producir dos
líneas, una en el segundo 10 y otra en el 90.
\end{ejercicio}

\section{Nivel 3: el puente}

\begin{ejercicio}[3.1 Una propiedad]
Añade a \cod{Player} la propiedad \cod{restantes} (cuántas canciones quedan en la cola) y
muéstrala en la pestaña Cola. Sigue la receta de la capítulo~\ref{cap:recetas}.
\end{ejercicio}

\begin{ejercicio}[3.2 Una función invocable]
Añade a \cod{Player} una función \cod{limpiar_cola()} que deje solo la pista actual. Llámala
desde una entrada nueva del menú contextual. ¿Qué pasa con \cod{index}? ¿Y con la precarga?
\end{ejercicio}

\begin{ejercicio}[3.3 Un singleton nuevo]
Crea un puente \cod{Notas} con una propiedad de texto y una función \cod{guardar}, que escriba
el texto en un archivo dentro de \cod{<config>/musicplayer/}. Regístralo en \archivo{build.rs}
y \archivo{main.rs}. Cuida lo siguiente: escribir en disco no debe bloquear el hilo de Qt.
\end{ejercicio}

\section{Nivel 4: hilos y red}

\begin{ejercicio}[4.1 Letra a demanda]
Añade a \cod{Library} una función que pida la letra de una canción cualquiera
(\cod{lyrics::fetch}) en otro hilo y la publique en una propiedad. Sigue el patrón de
\cod{search} de \archivo{library.rs}, \emph{con su contador} para descartar respuestas viejas.
¿Qué ocurre sin el contador si llamas dos veces seguidas?
\end{ejercicio}

\begin{ejercicio}[4.2 Medir una petición]
Mide, con tiempos reales, cuánto tarda \cod{ytm::browse("FEmusic_home")} (con
\cod{std::time::Instant} en una sonda nueva). Repítelo cinco veces: ¿qué dispersión hay?
\end{ejercicio}

\section{Nivel 5: la interfaz a fondo}

\begin{ejercicio}[5.1 Mover al final]
Añade «Mover al final» a las canciones de la cola en el menú de clic derecho (receta de
la capítulo~\ref{cap:recetas}). Pruébalo con \cod{MUSICPLAYER_DEBUG_CLICK} haciendo clic
derecho en una fila.
\end{ejercicio}

\begin{ejercicio}[5.2 Teclado en la cola]
En \archivo{QueueList.qml}, haz que las filas se puedan enfocar con Tab y que \texttt{Alt+$\uparrow$} y
\texttt{Alt+$\downarrow$} las muevan con \cod{Player.moveInQueue}. Respeta la norma: foco visible.
\end{ejercicio}

\begin{ejercicio}[5.3 Una skin completa]
Crea una skin en carpeta (\cod{skin.json} más una imagen de fondo con \cod{fondoOpacidad} y
\cod{fondoVelo}) en tu carpeta de skins, y elígela en Opciones. Haz una captura de Inicio, otra de
«Sonando» y otra de Opciones. ¿Se lee el texto sobre la imagen?
\end{ejercicio}

\begin{ejercicio}[5.4 Una mascota]
Crea una mascota con tres poses (\cod{reposo}, \cod{sonando}, \cod{clic}) y comprueba que
\cod{megusta} usa \cod{clic} al no existir. ¿Cuál es la cadena de sustitución completa para
\cod{playlist}?
\end{ejercicio}

\section{Nivel 6: rendimiento y diseño}

\begin{ejercicio}[6.1 Medir tu cambio]
Elige uno de los ejercicios anteriores que añada algo visible y mide su coste en CPU y memoria
con el método del capítulo~\ref{cap:rendimiento}: dos pasadas antes y dos después, en release,
con \cod{MUSICPLAYER_DEBUG_AUDIO}. Escribe la entrada de «Medidas» como si fuera a una ficha.
\end{ejercicio}

\begin{ejercicio}[6.2 Una animación que se detenga]
Añade un indicador que pulse (un círculo que cambia de opacidad) junto al título cuando suena
una canción. Debe \emph{pararse} en pausa, con la ventana oculta y con la opción
«Desenfoque» desactivada. Mide su coste. ¿Compensa un \cod{Timer} a 10 Hz frente a una
\cod{NumberAnimation}?
\end{ejercicio}

\begin{ejercicio}[6.3 Proponer alternativas]
La norma dice que el diseño visual se propone con \emph{dos a cuatro alternativas} con su
contrapartida. Propón tres formas de mostrar el tiempo restante (ejercicio 1.1): una
simple, una que cambie al pulsar y una que se anime. Escribe, para cada una, qué cuesta en
memoria, CPU y complejidad.
\end{ejercicio}

\begin{ejercicio}[6.4 Una ADR]
Escribe una ADR (\archivo{docs/runbooks/registrar-decision.md}) para una decisión ficticia:
«El tiempo restante se calcula en Rust y no en QML». Rellena \emph{alternativas descartadas} y
\emph{consecuencias} con argumentos reales. Ejecuta \cod{python3 scripts/docs.py check}.
\end{ejercicio}

\section{Pistas}

\begin{description}[style=nextline, leftmargin=0pt]
  \item[1.1] \cod{Theme.mmss(Player.duration - Player.position)}, en la línea que muestra la
  duración. Cuidado con que la duración puede ser 0 al arrancar.
  \item[2.1] Copia el caso \cod{"search"} de \cod{probe} en \archivo{main.rs}. Las zonas
  tienen título; busca la que se llame «Artistas» y toma su primer \cod{Item}.
  \item[2.2] \cod{seconds} es privada: la prueba va dentro de su propio archivo, en
  \cod{mod tests} con \cod{use super::*}.
  \item[3.1] Actualízala en \cod{set_tracks} y en \cod{load}; ¿y al quitar o mover pistas?
  \item[3.3] Escribe en un hilo (\cod{thread::spawn}) y publica el resultado con
  \cod{qt_thread().queue}. Usa \cod{write_atomic} de \archivo{settings.rs} como modelo.
  \item[4.1] Busca \cod{search_request} en \archivo{library.rs}.
  \item[5.2] \cod{activeFocusOnTab: true} y \cod{Keys.onPressed} mirando
  \cod{event.modifiers & Qt.AltModifier}.
  \item[6.2] Mira cómo \cod{Ambient.drifting} combina \cod{visible}, \cod{windowShown} y
  \cod{Player.playing}.
\end{description}

\section{Soluciones resumidas}

\subsection*{1.1}

\begin{qml}
Text {
    text: "-" + Theme.mmss(Math.max(0, Player.duration - Player.position))
}
\end{qml}

Solo cambia el texto; la duración es la misma propiedad que ya se leía.

\subsection*{2.1}

\begin{rust}
"artist" => {
    let zones = ytm::search(argument);
    if let Ok(zones) = &zones {
        let first = zones.iter().find(|z| z.titulo == "Artistas").and_then(|z| z.items.first());
        match first {
            Some(item) => println!("{} ({})", item.title, item.browse_id),
            None => println!("sin artistas"),
        }
    }
    return true;
}
\end{rust}

\subsection*{2.2}

\begin{rust}
#[test]
fn segundos_de_un_reloj() {
    assert_eq!(seconds("3:27"), 207.0);
    assert_eq!(seconds("1:02:03"), 3723.0);
    assert_eq!(seconds(""), 0.0);
}
\end{rust}

\subsection*{3.1}

La propiedad se declara en el puente y se inicializa en \cod{Default}. Se actualiza en cada
sitio donde cambian la cola o el índice: \cod{set_tracks}, \cod{load}, \cod{remove_from_queue}
y \cod{move_in_queue} (aunque mover no cambia el total, sí el \emph{índice actual}). Si la
calculas siempre en un único método auxiliar, no se te olvida ningún sitio.

\subsection*{4.1}

Un esbozo (los nombres de \cod{letra_request} y \cod{letra} son inventados; añádelos al puente):

\begin{rust}
fn letra_de(mut self: Pin<&mut Self>, titulo: &QString, artista: &QString) {
    let request = self.letra_request + 1;
    self.as_mut().rust_mut().letra_request = request;
    let query = lyrics::Query::new(&titulo.to_string(), &artista.to_string(), 0.0);
    let qt = self.qt_thread();
    thread::spawn(move || {
        let found = lyrics::fetch(lyrics::Source::Lrclib, &query);
        let _ = qt.queue(move |mut library| {
            if library.letra_request != request { return; }   // llegó tarde
            let json = serde_json::to_string(&found.map(|f| f.lines)).unwrap_or_default();
            library.as_mut().set_letra(QString::from(&json));
        });
    });
}
\end{rust}

Sin el contador, la primera respuesta (más lenta) puede llegar después de la segunda y
\emph{sobrescribirla} con una letra que ya no interesa.

\subsection*{6.2}

Con \cod{running: root.visible && Player.playing && Theme.blur && windowShown} en un
\cod{Timer} (o en la animación). Con un \cod{Timer} a 10 Hz, el coste es casi nulo; con una
\cod{NumberAnimation} a 60 Hz, cada fotograma redibuja la ventana y el vidrio desenfocado
detrás de ella: es exactamente el caso de la deriva del ambiente.
